\documentclass{article}
\usepackage{pathfinder-paper}

\title{What a Reviewed Quote Can Establish About Double-Auction Welfare}
\pathfinderpapers{2609.02580}{Competitive Market Behavior of LLMs}{2608.18167}{Adversarial Review:
Structured Disagreement for Grounded Agentic Code Review}

\begin{document}
\maketitle

\begin{abstract}
We examine what evidence-typed reviewer--critic scrutiny of an LLM trader's proposed quote could establish
in the double auction of Struski et al.\ and what it could not. Under that paper's stated values, maximum
surplus is \(7.50\) at either five or six trades, whereas the best seven-trade allocation yields \(7.00\).
Seven individually profitable matches remain possible. Thus a review that increases crossings or trade
count need not improve allocative efficiency. An exact split of surplus regret into quantity and selection
contributions gives an evaluation measure, not a causal explanation. The underlying price-cancellation and
inefficient-matching observations are established in earlier auction work; the contribution here is their
application to this pair and a bounded proposal for testing the unmeasured review intervention.
% ledger: 4, 5, 9, 13, 17, 20, 22
\end{abstract}

\section{Established setting and question}

Struski et al.\ study LLM traders in a continuous double auction with eleven single-unit buyers and eleven
single-unit sellers, each side having reservation values from \(0.75\) to \(3.25\) in \(0.25\) steps. The
full schedules imply a competitive price of \(2.00\) and the paper's equilibrium quantity of six. A round
selects one untraded agent at a time and stops after at most 300 iterations. A crossing order trades
immediately at the resting quote. Traders know their own reservation value, standing quotes, and
histories, and are instructed to maximise private profit. The authors report incomplete price convergence
and allocative efficiency, including rounds spent on one-cent quote improvements without completing
profitable trades. They already measure allocative efficiency; our argument concerns how to interpret
trade count and a proposed intervention, not a correction of their reported efficiency figures
\cite{struski2026}.
% ledger: 1, 2, 5, 18, 24

Qiu and Gill's Adversarial Review (AR) freezes a coding artefact while a reviewer and critic scrutinise a
review before the main agent edits it. In their SWE-PRBench review task, a text constraint distinguishing
evidence from epistemic concern raises review \(F_1\) from \(0.457\) for naive AR to \(0.533\). Their
SWE-bench Verified AR experiment instead uses the baseline two-verdict protocol. Neither study tests
auction quotes or auction welfare \cite{qiu2026}. The paired question is therefore whether such review
changes realised surplus when a trader's proposed quote is the frozen artefact. The sign of that effect is
unknown.
% ledger: 1, 10, 17, 18

\section{Accounting result for Q's values}

Let \(W\) denote realised buyer profit plus seller profit in one round, and \(q\) its number of trades.
For a completed match with buyer value \(v\), seller cost \(c\), and transaction price \(p\),
\[
(v-p)+(p-c)=v-c.
\]
Changing only the price for the same completed buyer--seller matches changes the division of rent but
leaves \(W\) fixed. Preventing a match or changing the participating traders can change \(W\). The
price-cancellation observation itself is standard and has appeared in prior continuous-double-auction
analysis \cite{wah2017}.
% ledger: 4, 18

Write \(v_{(k)}\) for the \(k\)-th highest buyer value and \(c_{(k)}\) for the \(k\)-th lowest seller
cost. For an allocation with exactly \(q\) trades, the largest possible surplus under Q's stated schedule
is
\[
W_q=\sum_{k=1}^{q}\bigl(v_{(k)}-c_{(k)}\bigr),
\qquad W^\star=\max_{0\leq q\leq11} W_q.
\]
For any realised allocation with \(q\) trades, the following is an exact ex post identity:
\[
W^\star-W=\underbrace{(W^\star-W_q)}_{\text{quantity contribution}}
            +\underbrace{(W_q-W)}_{\text{selection contribution}}.
\]
The second term measures the loss from which buyers and sellers traded at that count; it does not depend
on how those traders were paired once they are selected. The split is an accounting device, not a causal
mediation result.
% ledger: 5, 13, 20, 22

\begin{proposition}
Under Q's stated reservation values, \(W^\star=W_5=W_6=7.50\) and \(W_7=7.00\). There is a feasible
seven-trade matching in which every match has strictly positive surplus. Consequently, neither a move from
five to six trades nor a move from six to seven trades, by itself, establishes a welfare improvement.
% ledger: 5, 9, 11, 13
\end{proposition}

\begin{proof}
The successive differences \(v_{(k)}-c_{(k)}\) begin \(2.50,2.00,1.50,1.00,0.50,0,-0.50\), yielding the
stated maxima. For seven trades, match buyers with values \(3.25,3.00,\ldots,1.75\) to sellers with costs
\(2.25,2.00,\ldots,0.75\), respectively. Each pair has gain \(1.00\) and can transact at a price strictly
between its own value and cost. These seven gains sum to \(7.00\). The first two conclusions compare the
best allocations at each count; an actual six-trade allocation can still be worse than an actual
five-trade allocation if selection changes.
% ledger: 5, 9, 11, 13
\end{proof}

The proposition rebuts the initial inference that Q's reported mean above six trades must indicate an
error or a violation of trader instructions. Heterogeneous transaction prices permit seven individually
profitable bilateral matches. Q also states that its exchange does not enforce reservation bounds,
although its trader prompts instruct agents to respect them. We do not infer from the reported mean
whether any instruction was violated.
% ledger: 7, 9, 11, 24

\section{Prior work and interpretation}

The main inefficient-matching step is already published. Wah, Wright and Wellman give a
continuous-double-auction example in which two profitable matches yield less surplus than a different
single match, and explain that the transaction price cancels from the two traders' combined surplus
\cite{wah2017}. We claim no novelty for that principle or for the algebraic regret identity. What is
specific here is the zero-surplus sixth margin in Q's schedule, the feasible profitable seventh trade, and
the implication that a P-style quote review cannot be judged by crossings or volume alone. These are
analytic statements about Q's design, not observations of a reviewed market.
% ledger: 4, 5, 9, 12, 13, 18, 20

Adjacent LLM work already uses a critic to shape bidding rewards in an electricity continuous double
auction \cite{jadhav2025}, and checks protocol invariants in negotiation and auction messages
\cite{albayaydh2026}. Debate in simulated portfolio allocation has also been studied \cite{huang2026}.
These do not supply the missing outcome for a frozen-quote reviewer--critic intervention in Q's market.
The reviewer's evidence can establish current facts such as whether a resting ask gives a buyer known
immediate profit; it cannot establish that crossing maximises expected private profit without a model of
later fills. In a simplified choice between crossing ask \(a\) now and posting a lower bid \(b\) that
later fills with probability \(r\), crossing is privately preferable exactly when \(v-a>r(v-b)\). Q does
not identify \(r\). The analyst's full value schedule and equilibrium price must therefore remain outside
trader and reviewer prompts.
% ledger: 6, 8, 12, 14, 15, 18

\section{Open question and test}

A direct test would randomise whole market runs between Q's unreviewed baseline, a token-matched
extra-deliberation control, and evidence-typed reviewer--critic review. It would keep Q's private-profit
objective, trader information, order book, and iteration cap fixed, and save each proposed and final
action. The primary outcome would be realised \(W/W^\star\); secondary measures would include the two
accounting contributions, trader profits, crossings, and calls, tokens, and wall time. Randomisation can
estimate effects on these outcomes, but the accounting split alone cannot show that an effect was caused
through trade count rather than trader selection.
% ledger: 15, 17, 20, 22, 24

\section{Limitations}

No such intervention has been run here, so this paper makes no claim that review raises, lowers, or
preserves surplus. The calculation uses Q's stated symmetric value schedule and single-unit design; it
does not establish the same numerical margins for other schedules or models. The simple fill-probability
inequality omits competing arrivals and strategic responses. Q's sequential mechanism can hold the book
fixed during synchronous review, making added calls a resource cost in that setting; stale quotes would
require an asynchronous market experiment. Whether any welfare effect survives other models or value
schedules, and how forecast objections could be calibrated, remain open.
% ledger: 3, 8, 15, 16, 17, 20, 24

\clearpage
\bibliographystyle{plainurl}
\bibliography{references}
\end{document}
